
\section{Depth-Two Misadjustment Bound}

The remaining stationary non-martingale contribution is
\begin{equation}
\begin{aligned}
R_n^{\text{mis, RR}} \coloneqq \frac{1}{\sqrt{n}} \sum_{k = 0}^{n - 1}
                    (2 R_k^{(\alpha)} - R_k^{(2 \alpha)}),
\end{aligned}
\label{eq:R-mis-def}
\end{equation}
where \(R_k^{(\alpha)} \coloneqq J_k^{(1, \alpha)} + H_k^{(1, \alpha)}\).
% A direct kernel-difference route only gives
% \(O(\sqrt{n} \, \alpha)\), which is too crude at
% \(\alpha \asymp n^{-1 / 2}\). The depth-two route below uses
% Lemma~\ref{lem:levin-prop-2}, Lemma~\ref{lem:levin-cor-6}, Lemma~\ref{lem:levin-prop-8}, and
% Lemma~\ref{lem:levin-prop-9} to recover the target
% \(n^{-1 / 4} \, \text{polylog}(n)\) rate.

\textbf{Depth-two refinement.} For \(\ell \ge 1\) define
\begin{equation}
\begin{aligned}
J_n^{(\ell, \alpha)} \coloneqq (I - \alpha \overline{A}) \, J_{n - 1}^{(\ell, \alpha)}
                       - \alpha \widetilde{A}(Z_n) \, J_{n - 1}^{(\ell - 1, \alpha)},
\end{aligned}
\end{equation}
\begin{equation}
\begin{aligned}
H_n^{(\ell, \alpha)} \coloneqq (I - \alpha A(Z_n)) \, H_{n - 1}^{(\ell, \alpha)}
                       - \alpha \widetilde{A}(Z_n) \, J_{n - 1}^{(\ell, \alpha)},
\end{aligned}
\end{equation}
with \(J_0^{(\ell, \alpha)} = H_0^{(\ell, \alpha)} = 0\) and the \(\ell = 0\)
processes as in Section~\ref{sec:last_iterate}. Then
\begin{equation}
\begin{aligned}
H_n^{(\ell, \alpha)} = J_n^{(\ell + 1, \alpha)} + H_n^{(\ell + 1, \alpha)},
\quad \ell \ge 0.
\end{aligned}
\label{eq:depth-recursion}
\end{equation}
Applying \eqref{eq:depth-recursion} with \(\ell = 1\) refines \(R_k^{(\alpha)}\) to
\begin{equation}
\begin{aligned}
R_k^{(\alpha)} = J_k^{(1, \alpha)} + J_k^{(2, \alpha)} + H_k^{(2, \alpha)},
\end{aligned}
\end{equation}
so
\begin{equation}
\begin{aligned}
R_n^{\text{mis, RR}} = T_n^{(1)} + T_n^{(2)} + T_n^{(H)},
\end{aligned}
\label{eq:R-mis-split}
\end{equation}
\begin{equation}
\begin{aligned}
T_n^{(j)} \coloneqq \frac{1}{\sqrt{n}} \sum_{k = 0}^{n - 1}
              \left(2 J_k^{(j, \alpha)} - J_k^{(j, 2 \alpha)}\right),
\quad j \in \left\{1, 2\right\},
\end{aligned}
\end{equation}
with \(T_n^{(H)}\) defined identically with \(H^{(2)}\) in place of \(J^{(j)}\).
Each piece is bounded separately.
% The direct Levin working forms used in this section are collected in
% Section~\ref{sec:external-direct-inputs}.

\textbf{Finite-past truncation.} Fix \(q \ge 2\), \(2 \le p \le q / 2\), and
\(0 < w \le \alpha_{\text{stat}}(q)\). Let \((Z_t)_{t \in \mathbb{Z}}\) be a two-sided stationary
version of the driving Markov chain, and write \(B_w \coloneqq I - w \overline{A}\).
For \(m \ge 1\) start at time \(-m\) from zero,
\begin{equation}
\begin{aligned}
J_{-m,m}^{(0,w)} = J_{-m,m}^{(1,w)} = J_{-m,m}^{(2,w)}
  = H_{-m,m}^{(2,w)} = 0,
\end{aligned}
\end{equation}
and evolve, for \(r = -m + 1, \dots, 0\),
\begin{equation}
\begin{aligned}
J_{r,m}^{(0,w)}
  = B_w J_{r - 1,m}^{(0,w)} - w \varepsilon(Z_r),
\end{aligned}
\end{equation}
\begin{equation}
\begin{aligned}
J_{r,m}^{(\ell,w)}
  = B_w J_{r - 1,m}^{(\ell,w)}
    - w \widetilde{A}(Z_r) J_{r - 1,m}^{(\ell - 1,w)},
\quad \ell \in \left\{1, 2\right\},
\end{aligned}
\end{equation}
\begin{equation}
\begin{aligned}
H_{r,m}^{(2,w)}
  = (I - w A(Z_r)) H_{r - 1,m}^{(2,w)}
    - w \widetilde{A}(Z_r) J_{r - 1,m}^{(2,w)}.
\end{aligned}
\end{equation}
\begin{lemma}\label{lem:finite-past-full-augmented-state}
  \textbf{(Stationary full depth-two augmented state.)}
  There exist constants \(C_{\text{fp}}, c_{\text{fp}} > 0\), depending only on the fixed
  problem parameters allowed by the chapter convention, on \(C_W\), \(D_J\), \(D_H\)
  from Lemma~\ref{lem:levin-prop-5-component}, Lemma~\ref{lem:levin-prop-8}, Lemma~\ref{lem:levin-prop-9}, and on
  \(C_{\text{prod}}\), \(c_{\text{prod}}\) from Lemma~\ref{lem:burn-product-stability}, such that, with
\begin{equation}
\begin{aligned}
  A_{\text{fp}}(p,q,w)
    \coloneqq C_{\text{fp}} (1 + d^{1 / q}) p^8 t_{\text{mix}}^5
       \frac{1}{a} \sqrt{w / a} \log^3(1 / (w a)),
\end{aligned}
\end{equation}
  the terminal vector of the finite-past truncation above,
\begin{equation}
\begin{aligned}
  (J_{0,m}^{(0,w)}, J_{0,m}^{(1,w)}, J_{0,m}^{(2,w)},
   H_{0,m}^{(2,w)})
\end{aligned}
\end{equation}
  is Cauchy in \(L_p\). More precisely, for \(m' \ge m \ge 1\),
\begin{equation}
\begin{aligned}
  \lVert  H_{0,m'}^{(2,w)} - H_{0,m}^{(2,w)}  \rVert_{L_p}
    \le A_{\text{fp}}(p,q,w) \exp(-c_{\text{fp}} w a m / p),
\end{aligned}
\label{eq:finite-past-H2-cauchy}
\end{equation}
  The \(J\) coordinates converge to the invariant law
  \(\Pi_{J,2,w}\) of Lemma~\ref{lem:levin-invariant-depth-two-law}, and the \(L_p\) limit
\begin{equation}
\begin{aligned}
H_0^{(2,w)} \coloneqq \lim_{m \to \infty} H_{0,m}^{(2,w)}
\end{aligned}
\end{equation}
  satisfies the one-trajectory bound of Lemma~\ref{lem:levin-prop-9}.
\end{lemma}

\textit{Proof.} First control the \(J\)-coordinates. Compare the truncation started at
\(-m'\) with the one started at \(-m\) by applying the component contraction
\eqref{eq:levin-depth-two-component-contraction} after time \(-m\). The contraction
constant is \(C_W\) and its exponential rate is \(1 / 12\). The initial
depth-two cost at time \(-m\) is bounded by the elementary moment bounds for
\(J^{(0,w)}\) and \(J^{(1,w)}\), together with Lemma~\ref{lem:levin-prop-8} for
\(J^{(2,w)}\). Write
\begin{equation}
\begin{aligned}
Y_{r,m}^{(w)} \coloneqq (Z_{r + 1}, J_{r,m}^{(0,w)}, J_{r,m}^{(1,w)}, J_{r,m}^{(2,w)}).
\end{aligned}
\end{equation}
Thus there is a constant \(C_{\text{cost}}\), depending only on the
fixed problem parameters and on \(D_J\), such that
\begin{equation}
\begin{aligned}
\lVert  c_{J,2}^{(w)}(Y_{-m,m'}^{(w)}, Y_{-m,m}^{(w)})  \rVert_{L_p}
  \le C_{\text{cost}} (1 + d^{1 / q}) p^{7 / 2} t_{\text{mix}}^{5 / 2}
     \sqrt{w / a} \log^{3 / 2}(1 / (w a)).
\end{aligned}
\label{eq:finite-past-cost-bound}
\end{equation}
Consequently, for \(r \ge -m\) and \(\ell \in \left\{0,1,2\right\}\),
\begin{equation}
\begin{aligned}
\lVert  J_{r,m'}^{(\ell,w)} - J_{r,m}^{(\ell,w)}  \rVert_{L_p}
  \le A_J(p,q,w) \exp(-w a (r + m) / (12 p)),
\end{aligned}
\label{eq:finite-past-J-cauchy}
\end{equation}
where
\begin{equation}
\begin{aligned}
A_J(p,q,w)
  \coloneqq C_W C_{\text{cost}} (1 + d^{1 / q}) p^7 t_{\text{mix}}^5
     \sqrt{w / a} \log^3(1 / (w a)).
\end{aligned}
\label{eq:finite-past-AJ}
\end{equation}

Now subtract the two \(H^{(2,w)}\) finite-past recursions. With
\(\Delta H_r \coloneqq H_{r,m'}^{(2,w)} - H_{r,m}^{(2,w)}\) and similarly for
\(\Delta J_r^{(2,w)}\),
\begin{equation}
\begin{aligned}
\Delta H_0
  = \Gamma_{-m + 1:0}^{(w)} H_{-m,m'}^{(2,w)}
    - w \sum_{l = -m + 1}^0 \Gamma_{l + 1:0}^{(w)}
        \widetilde{A}(Z_l) \Delta J_{l - 1}^{(2,w)}.
\end{aligned}
\label{eq:finite-past-H2-difference}
\end{equation}

For the initial term, use the conditional product-stability estimate
\eqref{eq:burn-product-stability-conditional} and Lemma~\ref{lem:levin-prop-9}:
\begin{equation}
\begin{aligned}
\lVert  \Gamma_{-m + 1:0}^{(w)} H_{-m,m'}^{(2,w)}  \rVert_{L_p}
  &\le C_{\text{prod}} \exp(-c_{\text{prod}} w a m / p)
       \lVert  H_{-m,m'}^{(2,w)}  \rVert_{L_p} \\
  &\le C_{\text{prod}} D_H d^{1 / q} t_{\text{mix}}^{5 / 2} p^{7 / 2}
       w^{3 / 2} \log^{3 / 2}(1 / (w a))
       \exp(-c_{\text{prod}} w a m / p).
\end{aligned}
\label{eq:finite-past-H-initial-term}
\end{equation}

For the convolution term, combine \eqref{eq:burn-product-stability-conditional},
\(\lVert \widetilde{A}(z) \rVert \le C_A\), and \eqref{eq:finite-past-J-cauchy}:
\begin{equation}
\begin{aligned}
&w \sum_{l = -m + 1}^0
  \lVert  \Gamma_{l + 1:0}^{(w)} \widetilde{A}(Z_l) \Delta J_{l - 1}^{(2,w)}  \rVert_{L_p} \\
&\quad \le w C_{\text{prod}} C_A A_J(p,q,w)
   \sum_{l = -m + 1}^0
      \exp(-c_{\text{prod}} w a (-l) / p)
      \exp(-w a (l - 1 + m) / (12 p)).
\end{aligned}
\label{eq:finite-past-H-convolution-raw}
\end{equation}
The elementary convolution bound
\begin{equation}
\begin{aligned}
\sum_{l = -m + 1}^0
      \exp(-c_{\text{prod}} w a (-l) / p)
      \exp(-w a (l - 1 + m) / (12 p))
  \le C_{\text{conv}} \frac{p}{w a}
     \exp(-c_{\text{conv}} w a m / p)
\end{aligned}
\label{eq:finite-past-exp-convolution}
\end{equation}
holds with \(C_{\text{conv}}, c_{\text{conv}} > 0\) depending only on \(c_{\text{prod}}\) and
the numerical rate \(1 / 12\). Therefore the convolution term is bounded by
\begin{equation}
\begin{aligned}
C_{\text{prod}} C_A C_{\text{conv}} \frac{p}{a} A_J(p,q,w)
  \exp(-c_{\text{conv}} w a m / p).
\end{aligned}
\label{eq:finite-past-H-convolution-bound}
\end{equation}

Choose
\begin{equation}
\begin{aligned}
c_{\text{fp}} \coloneqq \frac{1}{2} \min(c_{\text{prod}}, c_{\text{conv}})
\end{aligned}
\end{equation}
and choose \(C_{\text{fp}}\) large enough to dominate the constants in
\eqref{eq:finite-past-H-initial-term} and \eqref{eq:finite-past-H-convolution-bound} after
substituting \eqref{eq:finite-past-AJ}. Since \(w \le \alpha_\infty\) implies
\(w a \le 1 / 2\), the initial-term factor \(w^{3 / 2}\) is absorbed by
\(a^{-1} \sqrt{w / a}\), and the powers
\(p^{7 / 2} t_{\text{mix}}^{5 / 2} \log^{3 / 2}(1 / (w a))\) are
absorbed by the displayed \(p^8 t_{\text{mix}}^5 \log^3(1 / (w a))\) in
\(A_{\text{fp}}(p,q,w)\). This gives \eqref{eq:finite-past-H2-cauchy}.

The \(J\)-bounds and \eqref{eq:finite-past-H2-cauchy} show that the full terminal vector
is Cauchy in \(L_p\). Passing to the \(L_p\) limit gives the stationary full
augmented-chain law.
% The limits are measurable functions of the two-sided stationary driving
% chain. Passing to the limit in the finite-past recursions is justified in
% \(L_p\), stationarity follows from shift-covariance, and the limiting
% \(H^{(2,w)}\) bound follows from Lemma~\ref{lem:levin-prop-9} and Fatou's lemma.
\begin{equation}
\begin{aligned}
\square
\end{aligned}
\end{equation}

Shifting the finite-past construction before taking the limit defines a
two-sided stationary process
\begin{equation}
\begin{aligned}
(Z_{k + 1}, J_k^{(0,w)}, J_k^{(1,w)}, J_k^{(2,w)}, H_k^{(2,w)})_{k \in \mathbb{Z}}.
\end{aligned}
\end{equation}
It solves the full depth-two recursions, and its law is the stationary full
augmented-chain law used below.

\textbf{Joint RR stationary construction.} Apply
Lemma~\ref{lem:finite-past-full-augmented-state} at \(w = \alpha\) and \(w = 2 \alpha\) on the
same two-sided stationary base chain.
The joint limit
\begin{equation}
\begin{aligned}
\left(
  Z_{k + 1},
  J_k^{(0, \alpha)}, J_k^{(1, \alpha)}, J_k^{(2, \alpha)},
  H_k^{(2, \alpha)},
  J_k^{(0, 2 \alpha)}, J_k^{(1, 2 \alpha)}, J_k^{(2, 2 \alpha)},
  H_k^{(2, 2 \alpha)}
\right)_{k \in \mathbb{Z}}
\end{aligned}
\end{equation}
is the stationary RR augmented state used below.
% All RR differences below, such as
% \(2 J_k^{(j, \alpha)} - J_k^{(j, 2 \alpha)}\) and
% \(2 H_k^{(2, \alpha)} - H_k^{(2, 2 \alpha)}\), are read under this joint
% construction.

\textbf{Stationary augmented-chain convention.} The estimates below use the stationary
versions constructed above. A zero-start full average is not obtained by a
single terminal contraction: summing startup contractions yields
\begin{equation}
\begin{aligned}
\frac{1}{\sqrt{n}} \sum_{k \ge 0} \rho_\alpha^k
  \asymp \frac{1}{\sqrt{n} \, \alpha a}
\end{aligned}
\end{equation}
when \(\rho_\alpha \approx \exp(-c \alpha a)\). This term is deferred to the burn-in
chapter.
% The one-step shift in
% \((Z_{t + 1}, J_t^{(0, w)}, J_t^{(1, w)}, J_t^{(2, w)}, H_t^{(2, w)})\) is
% intentional: \(J_t\) is the perturbation state after the observation
% \(Z_{t + 1}\) has updated the recursion, while base-chain covariance formulas
% may still use \(Z_0\) by stationarity. A practical arbitrary-start theorem
% requires either burned-in weights \(Q_{l,n_0}^{(\alpha)}\) or a separate transfer
% lemma with the accumulated startup dependence.

For the shifted-to-unshifted first-order transfer we use the local inverse
ceiling \(\alpha_{\text{inv}}\) defined in \eqref{eq:alpha-inv}.

\begin{lemma}\label{lem:stationary-limit-J1}
  \textbf{(Stationary-limit transfer for the centered first-order iterate.)}
  Fix \(w \in (0, \alpha_\infty]\) and set \(B_w \coloneqq I - w \overline{A}\). Assume
  the stationary augmented-chain convention above, \textbf{UGE 1},
  \(\pi(\widetilde{A}) = 0\), and \(\lVert  \varepsilon  \rVert_\infty < \infty\). Let
  \((J_t^{(0, w)}, J_t^{(1, w)})_{t \in \mathbb{Z}}\) be the stationary two-sided
  solution of the first two perturbation recursions, and set
  \(T_t^{(1, w)} \coloneqq B_w \, J_t^{(1, w)}\). Then, with
\begin{equation}
\begin{aligned}
  \Phi_+(p, w) \coloneqq 1 + p^{3 / 2} \, t_{\text{mix}}^{1 / 2} / a
               + p^{1 / 2} \, t_{\text{mix}}^{3 / 2} \sqrt{w / a},
\end{aligned}
\end{equation}
  there exists a constant \(C_{\text{stat,1}}\), depending only on the constants in
  the zero-start centered shifted first-order bound in Lemma~\ref{lem:last-shifted-first-order},
  such that
  for every \(p \ge 2\) and every \(t \in \mathbb{Z}\),
\begin{equation}
\begin{aligned}
  \lVert  T_t^{(1, w)} - \mathbb{E}_\pi T_t^{(1, w)}  \rVert_{L_p}
    \le C_{\text{stat,1}} \, w \, \Phi_+(p, w).
\end{aligned}
\end{equation}
  Consequently, whenever \(w \le \alpha_{\text{inv}}\),
\begin{equation}
\begin{aligned}
  \lVert  J_t^{(1, w)} - \mathbb{E}_\pi J_t^{(1, w)}  \rVert_{L_p}
    \le 2 C_{\text{stat,1}} \, w \, \Phi_+(p, w).
\end{aligned}
\end{equation}
\end{lemma}

\textit{Proof.} Use finite-past zero-start versions
\(J_{t, m}^{(0, w)}\) and \(J_{t, m}^{(1, w)}\). By stationarity and
Lemma~\ref{lem:last-shifted-first-order},
\begin{equation}
\begin{aligned}
\lVert  B_w J_{t, m}^{(1, w)} - \mathbb{E} B_w J_{t, m}^{(1, w)}  \rVert_{L_p}
  \le C_{\text{stat,1}} \, w \, \Phi_+(p, w).
\end{aligned}
\end{equation}
The deterministic-product expansions give
\begin{equation}
\begin{aligned}
J_{t, m}^{(0, w)}
  = -w \sum_{r = 0}^{m - 1} B_w^r \varepsilon(Z_{t - r}),
\end{aligned}
\end{equation}
and the corresponding \(J_{t,m}^{(1,w)}\) series is Cauchy for fixed admissible
\(w\). Passing to the limit gives the bound for \(T_t^{(1,w)}\). The bound for
\(J_t^{(1,w)}\) follows from
\begin{equation}
\begin{aligned}
J_t^{(1,w)} - \mathbb{E}_\pi J_t^{(1,w)}
  = B_w^{-1}\bigl(T_t^{(1,w)} - \mathbb{E}_\pi T_t^{(1,w)}\bigr).
\end{aligned}
\end{equation}
\(\square\)

% The finite-past domination is fixed-\(w\); later estimates keep the displayed
% \(w\)-dependence through \(\Phi_+(p,w)\) and use constants independent of \(n\).

\textbf{Telescoping identity for \(J^{(1)}\).} Summing
\begin{equation}
\begin{aligned}
J_k^{(1,\alpha)}
  = (I-\alpha\overline{A})J_{k-1}^{(1,\alpha)}
    -\alpha\widetilde{A}(Z_k)J_{k-1}^{(0,\alpha)}
\end{aligned}
\end{equation}
from \(k = 1\) to \(n\) and rearranging gives, for an arbitrary initial value,
\begin{equation}
\begin{aligned}
\overline{A} \sum_{k = 0}^{n - 1} J_k^{(1, \alpha)}
  = -\sum_{k = 1}^n \widetilde{A}(Z_k) \, J_{k - 1}^{(0, \alpha)}
    + \frac{1}{\alpha} \, \left(J_0^{(1, \alpha)} - J_n^{(1, \alpha)}\right).
\end{aligned}
\end{equation}
In stationarity,
\begin{equation}
\begin{aligned}
\mathbb{E}_\pi\!\left[\widetilde{A}(Z_1)J_0^{(0,\alpha)}\right]
  = -\overline{A}\,\mathbb{E}_\pi\!\left[J_\infty^{(1,\alpha)}\right].
\end{aligned}
\end{equation}
% Under the stationary augmented-chain convention,
% \((J_{k - 1}^{(0, w)}, Z_k)\) has the same law as
% \((J_0^{(0, w)}, Z_1)\), so \(\overline{\psi}_w\) is centered for each summand.
Subtracting \(n \, \mathbb{E}_\pi [J_\infty^{(1, \alpha)}]\) from both sides and
applying \(\overline{A}^{-1}\) gives the \textbf{centered telescoping identity}
\begin{equation}
\begin{aligned}
\sum_{k = 0}^{n - 1} \left(J_k^{(1, \alpha)} - \mathbb{E}_\pi [J_\infty^{(1, \alpha)}]\right)
  = -\overline{A}^{-1} \sum_{k = 1}^n
      \overline{\psi}_\alpha (J_{k - 1}^{(0, \alpha)}, Z_k)
    + \frac{1}{\alpha} \, \overline{A}^{-1}
      \left(J_0^{(1, \alpha)} - J_n^{(1, \alpha)}\right).
\end{aligned}
\label{eq:J1-telescope}
\end{equation}
% This identity links the PR-average of \(J^{(1)}\) to the centered bilinear sum
% bounded in Lemma~\ref{lem:levin-cor-6}.

\begin{lemma}\label{lem:T1-bound}
  \textbf{(Stationary first-order misadjustment bound.)}
  Assume the stationary augmented-chain convention above, \textbf{UGE 1},
  \(\pi(\widetilde{A}) = 0\), \(\pi(\varepsilon) = 0\),
  \(\lVert  \varepsilon  \rVert_\infty < \infty\), and
  \(0 < \alpha\), \(2 \alpha \le \alpha_\infty\), and
  \(2 \alpha \le \alpha_{\text{inv}}\). Set
\begin{equation}
\begin{aligned}
  \Phi_+(p, \alpha) \coloneqq 1 + p^{3 / 2} \, t_{\text{mix}}^{1 / 2} / a
                   + p^{1 / 2} \, t_{\text{mix}}^{3 / 2} \sqrt{\alpha / a}.
\end{aligned}
\end{equation}
  There exists a constant \(C_{\text{mis,1}}\) depending on
  \(\lVert\overline{A}\rVert\), \(\lVert\overline{A}^{-1}\rVert\),
  \(\kappa_Q\), \(C_A\), \(\lVert\varepsilon\rVert_\infty\),
  \(t_{\text{mix}}\), and the universals \(c_{W,1}\), \(c_{W,2}\)
  of Lemma~\ref{lem:levin-cor-6}, and the constant of the centered bound in
  Lemma~\ref{lem:last-shifted-first-order}, such that for every \(p \ge 2\) and every
  \(n \ge 2\),
\begin{equation}
\begin{aligned}
  \lVert  T_n^{(1)}  \rVert_{L_p}
    &\le C_{\text{mis,1}} \sqrt{n} \, \alpha^2
     + C_{\text{mis,1}} \, p^{3 / 2} \sqrt{\alpha} \\
    &\quad + C_{\text{mis,1}} \, p^3 \, (\alpha n)^{-1 / 2} \, \log^{1 / 2}(1 / (\alpha a))
     + C_{\text{mis,1}} \, \Phi_+(p, \alpha) \, n^{-1 / 2}.
\end{aligned}
\end{equation}
\end{lemma}

\textit{Proof.} Decompose \(T_n^{(1)} = T_n^{\text{(1, b)}} + T_n^{\text{(1, c)}}\) via the
bias-fluctuation split
\begin{equation}
\begin{aligned}
T_n^{\text{(1, b)}} \coloneqq \sqrt{n} \, \left(
                    2 \, \mathbb{E}_\pi [J_\infty^{(1, \alpha)}]
                    - \mathbb{E}_\pi [J_\infty^{(1, 2 \alpha)}]
                  \right),
\end{aligned}
\end{equation}
\begin{equation}
\begin{aligned}
T_n^{\text{(1, c)}} \coloneqq \frac{1}{\sqrt{n}} \sum_{k = 0}^{n - 1} \sum_{w \in \left\{\alpha, 2 \alpha\right\}} c_w
                   \left(J_k^{(1, w)} - \mathbb{E}_\pi [J_\infty^{(1, w)}]\right),
\quad c_\alpha = 2, \, c_{2 \alpha} = -1.
\end{aligned}
\end{equation}
% The bias is deterministic; the centered piece carries the \(L_p\) fluctuation.

\textbf{Bias.} By Lemma~\ref{lem:levin-prop-2},
\(\mathbb{E}_\pi [J_\infty^{(1, w)}] = w \, \Delta + R(w)\) with
\(\lVert  R(w)  \rVert \le 12 \, \lVert  \overline{A}^{-1}  \rVert \, C_A^2 \, t_{\text{mix}}^2 \, w^2 \, \lVert  \varepsilon  \rVert_\infty\).
The leading \(w \, \Delta\) cancels, leaving
\begin{equation}
\begin{aligned}
\lVert  2 \, \mathbb{E}_\pi [J_\infty^{(1, \alpha)}] - \mathbb{E}_\pi [J_\infty^{(1, 2 \alpha)}]  \rVert
  \le 2 \, \lVert  R(\alpha)  \rVert + \lVert  R(2 \alpha)  \rVert
  \le 72 \, \lVert  \overline{A}^{-1}  \rVert \, C_A^2 \, t_{\text{mix}}^2 \, \alpha^2 \, \lVert  \varepsilon  \rVert_\infty.
\end{aligned}
\end{equation}
Multiplying by \(\sqrt{n}\) gives the first term.

\textbf{Centered.} Apply \eqref{eq:J1-telescope} at \(w \in \left\{\alpha, 2 \alpha\right\}\):
\begin{equation}
\begin{aligned}
\lVert  T_n^{\text{(1, c)}}  \rVert_{L_p}
  &\le \frac{1}{\sqrt{n}} \sum_{w \in \left\{\alpha, 2 \alpha\right\}} |c_w| \Biggl(
       \lVert  \overline{A}^{-1}  \rVert \cdot
       \lVert  \sum_{k = 1}^n \overline{\psi}_w (J_{k - 1}^{(0, w)}, Z_k)  \rVert_{L_p} \\
    &\quad \quad \quad \quad
       + \frac{1}{w} \, \lVert  \overline{A}^{-1}  \rVert \cdot
         \lVert  J_0^{(1, w)} - J_n^{(1, w)}  \rVert_{L_p}
     \Biggr).
\end{aligned}
\end{equation}
For the bilinear sum,
\begin{equation}
\begin{aligned}
\lVert  \sum_{k = 1}^n \overline{\psi}_w  \rVert_{L_p}
  \le c_{W, 1} \, p^{3 / 2} \sqrt{w n}
   + c_{W, 2} \, p^3 \, w^{-1 / 2} \, \log^{1 / 2}(1 / (w a)),
\end{aligned}
\end{equation}
so
\begin{equation}
\begin{aligned}
\frac{1}{\sqrt{n}} \lVert  \sum \overline{\psi}_w  \rVert_{L_p}
  \le \sqrt{2} c_{W, 1} \, p^{3 / 2} \sqrt{\alpha}
   + \sqrt{2} c_{W, 2} \, p^3 \, (\alpha n)^{-1 / 2} \, \log^{1 / 2}(1 / (\alpha a)).
\end{aligned}
\end{equation}
For the boundary term, Lemma~\ref{lem:stationary-limit-J1} gives
\begin{equation}
\begin{aligned}
\lVert  J_n^{(1, w)} - \mathbb{E} J_n^{(1, w)}  \rVert_{L_p} \le C w \, \Phi_+(p, w),
\end{aligned}
\end{equation}
and Lemma~\ref{lem:levin-prop-2} gives \(\lVert  \mathbb{E} J_n^{(1, w)}  \rVert \le C w\), hence
\begin{equation}
\begin{aligned}
\lVert  J_n^{(1, w)}  \rVert_{L_p} \le C w \, \Phi_+(p, w).
\end{aligned}
\end{equation}
The same bound holds for \(J_0^{(1, w)}\), so
\begin{equation}
\begin{aligned}
w^{-1}\lVert J_0^{(1,w)}-J_n^{(1,w)}\rVert_{L_p}
  \le C\Phi_+(p,w) \le \sqrt{2}\,C\Phi_+(p,\alpha)
\end{aligned}
\end{equation}
after absorbing
the factor \(2\) into \(C\).
Dividing by \(\sqrt{n}\) gives the last term. \(\square\)

\begin{lemma}\label{lem:T2H-bound}
  \textbf{(Stationary second-order and depth-two remainder bound.)}
  Under the assumptions of the previous lemma, for every \(p \ge 2\) and every
  \(q \ge 2\) satisfying \(p \le q / 2\) and
  \(2 \alpha \le \alpha_{\text{stat}}(q)\),
\begin{equation}
\begin{aligned}
  \lVert  T_n^{(2)}  \rVert_{L_p} + \lVert  T_n^{(H)}  \rVert_{L_p}
    \le C_{\text{mis,2}} \, (1 + d^{1 / q}) \, p^{7 / 2} \, t_{\text{mix}}^{5 / 2}
       \, \sqrt{n} \, \alpha^{3 / 2} \, \log^{3 / 2}(1 / (\alpha a)),
\end{aligned}
\end{equation}
  with \(C_{\text{mis,2}} \coloneqq 6 \, \sqrt{2}^3 \, (D_J + D_H)\).
\end{lemma}

\textit{Proof.} By the triangle inequality and Lemma~\ref{lem:levin-prop-8},
\begin{equation}
\begin{aligned}
\lVert  T_n^{(2)}  \rVert_{L_p}
  &\le \frac{1}{\sqrt{n}} \sum_{k = 0}^{n - 1} \sum_{w \in \left\{\alpha, 2 \alpha\right\}} |c_w| \, \lVert  J_k^{(2, w)}  \rVert_{L_p} \\
  &\le 3 \, \sqrt{n} \, \sup_{w \in \left\{\alpha, 2 \alpha\right\}} \lVert  J_k^{(2, w)}  \rVert_{L_p} \\
  &\le 3 \, D_J \, (2 \alpha)^{3 / 2}
      \, t_{\text{mix}}^{5 / 2} \, p^{7 / 2} \, \sqrt{n}
      \, \log^{3 / 2}(1 / (\alpha a)),
\end{aligned}
\end{equation}
The same argument with Lemma~\ref{lem:levin-prop-9} controls \(T_n^{(H)}\). \(\square\)
% Here \(|c_\alpha| + |c_{2 \alpha}| = 3\) absorbs the RR combination and
% \((2 \alpha)^{3 / 2}\) is the worse bound at the larger step.

\begin{theorem}\label{thm:misadjustment}
  \textbf{(Stationary PR-averaged RR misadjustment bound.)}
  Assume \textbf{UGE 1}, \(\pi(\widetilde{A}) = 0\), \(\pi(\varepsilon) = 0\),
  \(\lVert  \varepsilon  \rVert_\infty < \infty\), \(0 < \alpha\), and
  the stationary augmented-chain convention above. There exists a constant \(C\)
  depending on the universal and problem constants of the previous two lemmas
  such that for every
  \(p \ge 2\), every \(q \ge 2\) satisfying
  \(p \le q / 2\) and \(2 \alpha \le \alpha_{\text{stat}}(q)\), and every \(n \ge 2\),
\begin{equation}
\begin{aligned}
  \lVert  R_n^{\text{mis, RR}}  \rVert_{L_p}
    &\le C \sqrt{n} \, \alpha^2
     + C \, (1 + d^{1 / q}) \, p^{7 / 2} \, t_{\text{mix}}^{5 / 2}
         \, \sqrt{n} \, \alpha^{3 / 2} \, \log^{3 / 2}(1 / (\alpha a)) \\
    &\quad + C \, p^{3 / 2} \sqrt{\alpha}
     + C \, p^3 \, (\alpha n)^{-1 / 2} \, \log^{1 / 2}(1 / (\alpha a))
     + C \, \Phi_+(p, \alpha) \, n^{-1 / 2}.
\end{aligned}
\end{equation}
\end{theorem}

\textit{Proof.} From \eqref{eq:R-mis-split},
\(\lVert  R_n^{\text{mis, RR}}  \rVert_{L_p} \le \lVert  T_n^{(1)}  \rVert_{L_p} + \lVert  T_n^{(2)}  \rVert_{L_p} + \lVert  T_n^{(H)}  \rVert_{L_p}\).
The result follows from Lemma~\ref{lem:T1-bound} and Lemma~\ref{lem:T2H-bound}. \(\square\)

\begin{corollary}\label{cor:misadjustment-rate}
  \textbf{(Stationary balanced-scale misadjustment rate.)}
  At the working scale \(\alpha = c \, n^{-1 / 2}\), with
  \(p = \max(2, \left\lceil \log n \right\rceil)\) and \(q = \max(2 p, \left\lceil \log(e \, d) \right\rceil, 2)\), and for \(n\)
  large enough that \(2 \alpha \le \alpha_{\text{stat}}(q)\),
\begin{equation}
\begin{aligned}
\lVert  R_n^{\text{mis, RR}}  \rVert_{L_p} \le C \, \text{polylog}(n) \, n^{-1 / 4},
\end{aligned}
\end{equation}
  of the same polynomial order as the leading martingale Berry--Esseen rate,
  up to logarithmic factors.
\end{corollary}

\textit{Proof.} At \(\alpha = c \, n^{-1 / 2}\): \(\sqrt{n} \alpha^2 = c^2 \, n^{-1 / 2}\),
\(\sqrt{n} \alpha^{3 / 2} = c^{3 / 2} \, n^{-1 / 4}\),
\(\sqrt{\alpha} = c^{1 / 2} \, n^{-1 / 4}\),
\((\alpha n)^{-1 / 2} = c^{-1 / 2} \, n^{-1 / 4}\), and
\(\Phi_+(p, \alpha) \, n^{-1 / 2} = O(p^{3 / 2} \, n^{-1 / 2})\).
With \(p \asymp \log n\) and \(q \ge \log(e d)\), the displayed bound is
\(\text{polylog}(n) \, n^{-1 / 4}\). \(\square\)

This is a stationary \(n_0 = 0\) augmented-chain bound; the deterministic-start
version is handled in the burn-in transfer chapter.
